% Afinidades y coordenadas homogéneas (partición de 03-reflexiones-afinidades-homogeneas.tex, fase B)
\chapter{Afinidades y coordenadas homogéneas}\label{cap:afines}

% Números de los ejemplos y soluciones: apunte/scripts/cap06_afines.py

\section{La idea}

Las isometrías del capítulo~\ref{cap:reflexiones} conservan distancias; las similaridades conservan
ángulos y multiplican todas las distancias por un mismo factor. En este capítulo aflojamos
la rigidez un paso más: permitimos que una figura se estire distinto en cada dirección o se
incline, siempre que las rectas sigan siendo rectas y las paralelas sigan siendo paralelas.
Ese grupo de transformaciones es el de las \emph{afinidades}. Usamos «afinidad» y
«transformación afín» como sinónimos. Para escribirlas con una sola matriz introducimos las
\emph{coordenadas homogéneas}, que transforman la traslación en una multiplicación más.

\subsection{Afinidades: conservar rectas y paralelismo}\label{sec:afinidades-primera}

{\small\vspace{1mm}\linkUnidad{04_afines_homogeneas}{Simulación interactiva: Unidad 6 — Afinidades y coordenadas homogéneas}\par}

Una afinidad permite deformaciones más generales que un giro o una homotecia. Un cuadrado
puede convertirse en un paralelogramo y el círculo inscripto en una elipse inscripta; las
rectas que dividen al cuadrado en cuatro siguen dividiendo al paralelogramo en cuatro, y los
puntos medios siguen siendo puntos medios.

\begin{center}
\begin{tikzpicture}[scale=.86]
  % cuadrado con su círculo inscripto y sus medianas
  \begin{scope}[shift={(1,1)},scale=2.7]
    \draw[teal,line width=1.1pt] (0,0) rectangle (1,1);
    \draw[auxline] (.5,0)--(.5,1) (0,.5)--(1,.5);
    \draw[teal,line width=.8pt] (.5,.5) circle (.5);
  \end{scope}
  % su imagen afín: (0,0)->(7,1), (1,0)->(10.2,1.5), (0,1)->(5.9,3.7)
  \begin{scope}[cm={3.2,.5,-1.1,2.7,(7,1)}]
    \draw[rust,line width=1.1pt] (0,0)--(1,0)--(1,1)--(0,1)--cycle;
    \draw[auxline] (.5,0)--(.5,1) (0,.5)--(1,.5);
    \draw[rust,line width=.8pt] plot[domain=0:360,samples=72,smooth]
      ({.5+.5*cos(\x)},{.5+.5*sin(\x)});
  \end{scope}
  \draw[-{Latex[length=2.8mm]},tealdark,line width=1pt]
    (4.3,2.4)--(5.4,2.4) node[midway,above] {$f$};
  \node[labelback,text=teal] at (2.35,.55) {cuadrado y círculo};
  \node[labelback,text=rust] at (8.05,.55) {paralelogramo y elipse};
\end{tikzpicture}
\end{center}

\begin{idea}[Idea]
Una afinidad es una transformación lineal seguida de una traslación: conserva rectas,
paralelismo y razones sobre una recta, pero no distancias ni ángulos.
\end{idea}

En fórmulas, una afinidad del plano se escribe
\[
  f(x)=Ax+b,
\]
con \(A\) una matriz \(2\times2\) y \(b\) un vector. La parte \(A\) deforma la cuadrícula
(lleva el cuadrado unidad a un paralelogramo) y \(b\) desplaza el origen. Si \(\det A\neq0\),
ninguna dirección colapsa en un punto y la transformación se puede deshacer.

\subsection{Coordenadas homogéneas: puntos como rayos}\label{sec:homogeneas-primera}

{\small\vspace{1mm}\linkUnidad{04_afines_homogeneas}{Simulación interactiva: Unidad 6 — Afinidades y coordenadas homogéneas}\par}

\paragraph{Una motivación de campo.}
Un relevamiento con estación total se registra en un sistema local \((x,y)\): origen en la
estación, eje \(x\) hacia una referencia arbitraria. Para llevarlo al marco oficial
\((E,N)\) hay que girar un ángulo \(\theta\) (medido en sentido antihorario desde el eje
\(E\)) y trasladar al origen \((E_0,N_0)\):
\[
  \begin{pmatrix}E\\N\end{pmatrix}
  =
  \begin{pmatrix}\cos\theta&-\sin\theta\\ \sin\theta&\cos\theta\end{pmatrix}
  \begin{pmatrix}x\\y\end{pmatrix}
  +\begin{pmatrix}E_0\\N_0\end{pmatrix}.
\]
Son dos operaciones de naturaleza distinta: un producto y una suma. Si después hay que pasar
a otro marco (otra rotación y otra traslación), las cuentas se encadenan con sumas y
productos intercalados. Las coordenadas homogéneas agregan una tercera coordenada igual
a~\(1\) y escriben todo como un único producto por una matriz \(3\times3\):
\[
  \begin{pmatrix}E\\N\\1\end{pmatrix}
  =
  \begin{pmatrix}\cos\theta&-\sin\theta&E_0\\ \sin\theta&\cos\theta&N_0\\ 0&0&1\end{pmatrix}
  \begin{pmatrix}x\\y\\1\end{pmatrix}.
\]
Encadenar cambios de sistema es, entonces, multiplicar matrices \(3\times3\). La
georreferenciación de una imagen (ejemplo~\ref{ej:06-dron}) funciona igual.

\paragraph{El plano \(w=1\).}
Geométricamente, el plano afín se ve como el plano \(w=1\) dentro de un espacio
tridimensional de coordenadas \((x,y,w)\). El punto \((x,y)\) se representa por el vector
\((x,y,1)^T\). Cualquier múltiplo no nulo \(\mu\,(x,y,1)^T\) está sobre el mismo rayo que
parte del origen y corta al plano \(w=1\) en el mismo punto; por eso se dice que un punto es
un rayo.

\begin{center}
\begin{tikzpicture}[scale=.92]
  \coordinate (o) at (1,1);
  \draw[axisline,-{Latex[length=2.4mm]}] (o)--(8.2,1) node[right] {$y$};
  \draw[axisline,-{Latex[length=2.4mm]}] (o)--(-.2,-.1) node[below left] {$x$};
  \draw[axisline,-{Latex[length=2.4mm]}] (o)--(1,6) node[above] {$w$};
  \draw[teallight,fill=teallight!55,line width=.8pt]
    (-.2,3.2)--(6.8,3.2)--(8.2,4.25)--(1.2,4.25)--cycle;
  \node[labelback,text=tealdark] at (7.1,4.55) {plano \(w=1\)};
  \coordinate (p) at (4.8,3.74);
  \fill[rust] (p) circle (2.7pt) node[labelback,above] {\((x,y,1)\)};
  \draw[rust,line width=1.05pt,-{Latex[length=2.5mm]}] (o)--($(o)!1.38!(p)$);
  \fill[ink] (o) circle (2.4pt) node[labelback,below] {origen};
  \node[labelback,font=\small,text=rust] at (3.4,2.1)
    {un rayo representa un punto};
\end{tikzpicture}
\end{center}

\paragraph{Puntos en el infinito.}
Si un punto recorre una recta \(p+s\,d\) del plano \(w=1\), su rayo se va inclinando y, cuando
\(s\) crece, tiende a quedar horizontal, en la dirección \((d,0)\). Ese rayo horizontal no
corta al plano \(w=1\): no es un punto ordinario, sino el \emph{punto en el infinito} de la
recta. Todas las rectas paralelas a \(d\) tienden al mismo rayo; en ese sentido, las paralelas
se cortan en el infinito.

\begin{center}
\begin{tikzpicture}[scale=.9]
  \draw[axisline,-{Latex[length=2.2mm]}] (0,0)--(10.4,0) node[right] {\small dirección \(d\)};
  \draw[axisline,-{Latex[length=2.2mm]}] (0,0)--(0,2.6) node[above] {$w$};
  \draw[mainline] (-.3,1.6)--(10.2,1.6);
  \node[text=teal,font=\small,left] at (-.3,1.6) {\(w=1\)};
  \fill[ink] (0,0) circle (2.2pt) node[below] {\small origen};
  \foreach \s/\lab in {.8/{s=0},2.4/{s=1},4.8/{s=2},9.6/{s=4}}{
    \draw[rust,line width=.8pt] (0,0)--($(0,0)!1.1!(\s,1.6)$);
    \fill[rust] (\s,1.6) circle (2.2pt);
    \node[labelback,font=\scriptsize,below right] at (\s,1.55) {\(\lab\)};
  }
  \draw[gold,line width=1.1pt,-{Latex[length=2.5mm]}] (0,0)--(9.5,0.06);
  \node[labelback,font=\small,text=ink] at (6.6,.42) {rayo límite \((d,0)\): punto en el infinito};
\end{tikzpicture}
\end{center}

\begin{idea}[Idea]
Las coordenadas homogéneas agregan una coordenada igual a \(1\). Con ella, la traslación pasa
a ser una columna más de la matriz, componer afinidades es multiplicar matrices \(3\times3\),
y las direcciones aparecen como puntos con última coordenada \(0\).
\end{idea}

\section{Construcciones, cuentas y ejemplos}

Ahora formalizamos las afinidades, demostramos qué conservan, buscamos sus puntos fijos, las
escribimos en coordenadas homogéneas y las construimos a partir de datos geométricos: un eje
y una dirección, o tres puntos de control.

{\small\vspace{1mm}\linkUnidad{04_afines_homogeneas}{Simulación interactiva: Unidad 6 — Afinidades y coordenadas homogéneas}\par}

\subsection{Afinidades: idea geométrica y forma algebraica}\label{sec:afinidades-geom}

\begin{definicion}[Transformación afín]\label{def:06-afinidad}
Una \emph{transformación afín} o \emph{afinidad} del plano es una función
\(f\colon\R^2\to\R^2\) de la forma
\[
  \boxed{f(x)=Ax+b},
\]
donde \(A\) es una matriz \(2\times2\) con \(\det A\neq0\) y \(b\in\R^2\). La matriz \(A\) es
la \emph{parte lineal} y \(b\) la \emph{traslación}.
\end{definicion}

Las isometrías (\(A\) ortogonal) y las similaridades (\(A\) múltiplo de una ortogonal) del
capítulo~\ref{cap:reflexiones} son casos particulares. La figura muestra el efecto típico: un
paralelogramo va a otro paralelogramo.

\begin{center}
\begin{tikzpicture}[scale=.86]
  \coordinate (a) at (1,1);
  \coordinate (b) at (4.1,1.2);
  \coordinate (c) at (5.2,3.2);
  \coordinate (d) at (2.1,3);
  \coordinate (ap) at (7,1.2);
  \coordinate (bp) at (10.4,.7);
  \coordinate (cp) at (11.7,3.4);
  \coordinate (dp) at (8.3,3.9);
  \draw[teal,line width=1.05pt] (a)--(b)--(c)--(d)--cycle;
  \draw[rust,line width=1.05pt] (ap)--(bp)--(cp)--(dp)--cycle;
  \draw[-{Latex[length=2.8mm]},tealdark,line width=.9pt]
    (5.7,2.2)--(6.6,2.2) node[midway,above] {$f$};
  \node[font=\small,text=teal] at (3.1,.55) {paralelogramo};
  \node[font=\small,text=rust] at (9.5,.25) {paralelogramo};
\end{tikzpicture}
\end{center}

\begin{proposicion}[Qué conserva una afinidad]\label{prop:06-conserva}
Una afinidad \(f(x)=Ax+b\) lleva rectas en rectas, rectas paralelas en rectas paralelas, y
conserva la razón entre segmentos de una misma recta: si \(r=p+s_1d\) y \(q=p+s_2d\), entonces
\(f(r)-f(p)\) y \(f(q)-f(p)\) guardan la misma razón \(s_1:s_2\).
\end{proposicion}

\begin{demostracion}
La recta \(L=\{p+s\,d : s\in\R\}\), con \(d\neq0\), se transforma en
\[
  f(p+s\,d)=Ap+b+s\,Ad=f(p)+s\,Ad .
\]
Como \(\det A\neq0\), \(Ad\neq0\), y la imagen es la recta que pasa por \(f(p)\) con dirección
\(Ad\). Dos rectas paralelas tienen la misma dirección \(d\); sus imágenes tienen la misma
dirección \(Ad\) y siguen siendo paralelas. Finalmente, \(f(p+s_id)-f(p)=s_i\,Ad\) para
\(i=1,2\): los dos vectores son múltiplos de \(Ad\) con coeficientes \(s_1\) y \(s_2\), así que
la razón se conserva. En particular, el punto medio de un segmento va al punto medio de su
imagen.
\end{demostracion}

Las distancias y los ángulos, en cambio, no se conservan: el ejemplo siguiente los rompe a
propósito.

\begin{ejemplo}[Escalado no uniforme y corte]\label{ej:06-escalado}
El \emph{escalado no uniforme} \(S=\begin{pmatrix}2&0\\0&1\end{pmatrix}\) duplica las
abscisas y deja las ordenadas. El círculo unidad \(x^2+y^2=1\) va a la curva formada por los
puntos \((2x,y)\), es decir, a la elipse
\[
  \frac{X^2}{4}+Y^2=1 .
\]
Los puntos \((1,0)\), \((0,1)\) y \((0{,}707,\ 0{,}707)\) del círculo van a \((2,0)\), \((0,1)\) y
\((1{,}414,\ 0{,}707)\). Las distancias al origen dejan de ser iguales: la afinidad no es una
similaridad. El área pasa de \(\pi\approx3{,}142\) a \(2\pi\approx6{,}283\), porque
\(\det S=2\).

El \emph{corte} \(C=\begin{pmatrix}1&1\\0&1\end{pmatrix}\) desplaza cada punto horizontalmente
en una cantidad igual a su ordenada. El cuadrado unidad \((0,0),(1,0),(1,1),(0,1)\) va al
paralelogramo \((0,0),(1,0),(2,1),(1,1)\): los ángulos rectos dejan de serlo, pero
\(\det C=1\) y el área se conserva.
\end{ejemplo}

\begin{center}
\begin{tikzpicture}[scale=.95]
  \begin{scope}
    \draw[axisline,-{Latex[length=2mm]}] (-2.4,0)--(2.6,0) node[right] {$x$};
    \draw[axisline,-{Latex[length=2mm]}] (0,-1.4)--(0,1.6) node[above] {$y$};
    \draw[teal,line width=1pt] (0,0) circle (1);
    \draw[rust,line width=1pt] (0,0) ellipse (2 and 1);
    \fill[teal] (.707,.707) circle (2pt);
    \fill[rust] (1.414,.707) circle (2pt);
    \draw[auxline,-{Latex[length=1.8mm]}] (.75,.707)--(1.36,.707);
    \node[labelback,font=\small,text=teal] at (-.55,1.25) {círculo};
    \node[labelback,font=\small,text=rust] at (1.7,-1.15) {elipse \(S\)};
  \end{scope}
  \begin{scope}[xshift=6.2cm,yshift=-.6cm]
    \draw[axisline,-{Latex[length=2mm]}] (-.4,0)--(2.6,0) node[right] {$x$};
    \draw[axisline,-{Latex[length=2mm]}] (0,-.4)--(0,1.7) node[above] {$y$};
    \draw[teal,line width=1pt] (0,0) rectangle (1,1);
    \draw[rust,line width=1pt] (0,0)--(1,0)--(2,1)--(1,1)--cycle;
    \draw[auxline,-{Latex[length=1.8mm]}] (0,1.15)--(1,1.15);
    \node[labelback,font=\small,text=rust] at (1.9,1.4) {corte \(C\)};
  \end{scope}
\end{tikzpicture}
\end{center}

\subsection{Puntos fijos y forma afín}\label{sec:puntos-fijos-afin}

\begin{definicion}[Punto fijo]\label{def:06-punto-fijo}
Un punto \(p_\ast\) es \emph{fijo} para la afinidad \(f\) si \(f(p_\ast)=p_\ast\). Para
\(f(x)=Ax+b\) esto equivale al sistema lineal
\[
  Ap_\ast+b=p_\ast
  \quad\Longleftrightarrow\quad
  (I-A)\,p_\ast=b .
\]
\end{definicion}

Si \(\det(I-A)\neq0\) hay un único punto fijo; si \(\det(I-A)=0\) puede no haber ninguno
(una traslación) o haber una recta entera de puntos fijos (una reflexión). La tabla reúne las
transformaciones del capítulo~\ref{cap:reflexiones} escritas como afinidades; la fórmula de la
reflexión \(\sigma_L\) respecto de la recta \(L\colon n^Tx=h\) se deduce en ese capítulo.

\begin{center}
\begin{tabularx}{\textwidth}{>{\bfseries}l l X}
\toprule
Transformación & \(A,\ b\) & Puntos fijos\\
\midrule
Traslación \(x\mapsto x+b\) & \(A=I\) & ninguno si \(b\neq0\)\\
Rotación \(\rhoP{c,\theta}\) & \(A=R(\theta),\ b=(I-R(\theta))c\) & \(c\), si \(\theta\neq0\)\\
Homotecia \(\homot{o}{k}\) & \(A=kI,\ b=(1-k)o\) & \(o\), si \(k\neq1\)\\
Reflexión \(\sigma_L\), \(L\colon n^Tx=h\) &
\(A=I-\dfrac{2nn^T}{\lVert n\rVert^2},\
b=\dfrac{2h}{\lVert n\rVert^2}n\) &
todos los puntos de \(L\); \(b\) es normal a \(L\)\\
\bottomrule
\end{tabularx}
\end{center}

Si \(p_\ast\) existe, al tomarlo como nuevo origen mediante \(z=x-p_\ast\) la afinidad queda
puramente lineal:
\[
  f(x)-p_\ast=Ax+b-(Ap_\ast+b)=A(x-p_\ast)
  \qquad\Longrightarrow\qquad
  z\mapsto Az .
\]
Elegir bien el origen simplifica la transformación: alrededor de su punto fijo, toda afinidad
se ve como una transformación lineal. El ejemplo~\ref{ej:06-rot-escala} calcula un punto
fijo concreto.

\subsection{Coordenadas homogéneas}\label{sec:homogeneas}

\begin{definicion}[Coordenadas homogéneas]\label{def:06-homogeneas}
Las \emph{coordenadas homogéneas} del punto \((x,y)\) son cualquier terna
\((\mu x,\mu y,\mu)^T\) con \(\mu\neq0\); el representante normalizado es \((x,y,1)^T\). A la
inversa, una terna \((X,Y,W)^T\) con \(W\neq0\) representa el punto \((X/W,\ Y/W)\).
\end{definicion}

\begin{definicion}[Punto impropio]\label{def:06-impropio}
Una terna \((d_1,d_2,0)^T\neq0\) se llama \emph{punto impropio} o \emph{punto en el infinito}
de dirección \(d=(d_1,d_2)\). Todas las rectas paralelas a \(d\) comparten ese punto impropio.
\end{definicion}

Con esta escritura, la afinidad \(f(x)=Ax+b\), con \(A=\begin{pmatrix}a_{11}&a_{12}\\a_{21}&a_{22}\end{pmatrix}\)
y \(b=(b_1,b_2)\), se vuelve un único producto:
\[
\begin{pmatrix}x'\\y'\\1\end{pmatrix}
=
\begin{pmatrix}
a_{11}&a_{12}&b_1\\
a_{21}&a_{22}&b_2\\
0&0&1
\end{pmatrix}
\begin{pmatrix}x\\y\\1\end{pmatrix},
\qquad
M_f=\begin{pmatrix}A&b\\0\ \ 0&1\end{pmatrix}.
\]
El bloque superior izquierdo es \(A\), la última columna es \(b\) y la última fila es
\((0\ 0\ 1)\). Visto en \(\R^3\), \(M_f\) es una transformación \emph{lineal} que deja
invariante el plano \(w=1\): las afinidades del plano son exactamente las transformaciones
lineales invertibles de \(\R^3\) con última fila \((0\ 0\ 1)\). Sobre un punto impropio, la
traslación no actúa:
\[
  M_f\begin{pmatrix}d\\0\end{pmatrix}=\begin{pmatrix}Ad\\0\end{pmatrix}.
\]
Las direcciones sólo sienten la parte lineal, y un punto impropio va a otro punto impropio.
En el capítulo~\ref{cap:proyectivas} se permite que la última fila sea cualquiera; entonces los
puntos impropios pueden volverse ordinarios, y aparecen las transformaciones proyectivas.

\begin{proposicion}[Componer es multiplicar]\label{prop:06-composicion}
Si \(f(x)=A_1x+b_1\) y \(g(x)=A_2x+b_2\), la composición \(g\circ f\) (primero \(f\), después
\(g\)) es la afinidad
\[
  (g\circ f)(x)=A_2A_1\,x+(A_2b_1+b_2),
\]
y su matriz homogénea es el producto \(M_{g\circ f}=M_gM_f\).
\end{proposicion}

\begin{demostracion}
Sustituyendo, \(g(f(x))=A_2(A_1x+b_1)+b_2=A_2A_1x+A_2b_1+b_2\). El producto por bloques da lo
mismo:
\[
\begin{pmatrix}A_2&b_2\\0\ \ 0&1\end{pmatrix}
\begin{pmatrix}A_1&b_1\\0\ \ 0&1\end{pmatrix}
=
\begin{pmatrix}A_2A_1&A_2b_1+b_2\\0\ \ 0&1\end{pmatrix}.
\]
La matriz de la derecha actúa primero, como corresponde a vectores columna.
\end{demostracion}

\begin{ejemplo}[Georreferenciar una imagen de dron]\label{ej:06-dron}
Una ortofoto de \(400\times300\) píxeles tiene coordenadas de imagen \((u,v)\): origen en el
vértice superior izquierdo, \(u\) hacia la derecha y \(v\) hacia abajo. Cada píxel cubre
\(0{,}5\ \mathrm{m}\), y el vértice superior izquierdo cae en \(E=5000\), \(N=8000\). Entonces
\[
  x_E=5000+0{,}5\,u,\qquad y_N=8000-0{,}5\,v .
\]
El signo menos aparece porque \(v\) crece hacia abajo y \(N\) hacia arriba. En forma afín,
\[
  A=\begin{pmatrix}0{,}5&0\\0&-0{,}5\end{pmatrix},
  \qquad b=\begin{pmatrix}5000\\8000\end{pmatrix},
  \qquad \det A=-0{,}25 .
\]
El determinante es negativo: la transformación invierte la orientación. En la imagen, el
giro corto de \(u\) a \(v\) es horario; en el terreno, el de \(E\) a \(N\) es antihorario.
El módulo \(0{,}25\) dice que cada píxel (de área \(1\)) cubre \(0{,}25\ \mathrm{m}^2\).
En homogéneas,
\[
  \begin{pmatrix}x_E\\y_N\\1\end{pmatrix}
  =
  \begin{pmatrix}0{,}5&0&5000\\0&-0{,}5&8000\\0&0&1\end{pmatrix}
  \begin{pmatrix}u\\v\\1\end{pmatrix}.
\]
Para el píxel \((0,0)\) se obtiene \((5000,\ 8000)\); para el píxel \((200,100)\),
\[
  \begin{pmatrix}0{,}5\cdot200+5000\\-0{,}5\cdot100+8000\\1\end{pmatrix}
  =\begin{pmatrix}5100\\7950\\1\end{pmatrix}.
\]
Los otros vértices van a \((400,0)\mapsto(5200,\ 8000)\), \((0,300)\mapsto(5000,\ 7850)\) y
\((400,300)\mapsto(5200,\ 7850)\): la imagen cubre un rectángulo de
\(200\ \mathrm{m}\times150\ \mathrm{m}\). El punto impropio \((1,0,0)^T\) (dirección de las
filas de la imagen) va a \((0{,}5,\ 0,\ 0)^T\), dirección Este; la traslación no lo afecta.
\end{ejemplo}

\begin{center}
\begin{tikzpicture}[scale=.95]
  % imagen: 1 cm = 100 px, v hacia abajo
  \begin{scope}
    \draw[step=1,linegray,line width=.4pt] (0,-3) grid (4,0);
    \draw[teal,line width=1.1pt] (0,-3) rectangle (4,0);
    \draw[axisline,-{Latex[length=2mm]}] (0,0)--(4.6,0) node[right] {$u$};
    \draw[axisline,-{Latex[length=2mm]}] (0,0)--(0,-3.6) node[below] {$v$};
    \fill[ink] (0,0) circle (2.2pt) node[labelback,above left] {\small\((0,0)\)};
    \fill[rust] (2,-1) circle (2.4pt) node[labelback,above right] {\small\((200,100)\)};
    \draw[gold,line width=.9pt,-{Latex[length=2mm]}] (.7,0) arc (0:-90:.7);
    \node[labelback,font=\scriptsize] at (1.1,-.9) {horario};
    \node[labelback,font=\small,text=teal] at (2,-4.05) {imagen (píxeles)};
  \end{scope}
  \draw[-{Latex[length=2.8mm]},tealdark,line width=1pt] (5.2,-1.5)--(6.4,-1.5)
    node[midway,above] {$f$};
  % terreno: 1 cm = 50 m, origen del dibujo en (5000, 7850)
  \begin{scope}[xshift=7.6cm,yshift=-3cm]
    \draw[step=1,linegray,line width=.4pt] (0,0) grid (4,3);
    \draw[rust,line width=1.1pt] (0,0) rectangle (4,3);
    \draw[axisline,-{Latex[length=2mm]}] (-.3,0)--(4.6,0) node[right] {$E$};
    \draw[axisline,-{Latex[length=2mm]}] (0,-.3)--(0,3.6) node[above] {$N$};
    \fill[ink] (0,3) circle (2.2pt) node[labelback,left] {\small\((5000,8000)\)};
    \fill[rust] (2,2) circle (2.4pt) node[labelback,above right] {\small\((5100,7950)\)};
    \node[labelback,font=\scriptsize,below] at (0,-.3) {5000};
    \node[labelback,font=\scriptsize,below] at (4,-.3) {5200};
    \node[labelback,font=\scriptsize,left] at (-.3,0) {7850};
    \draw[gold,line width=.9pt,-{Latex[length=2mm]}] (.7,0) arc (0:90:.7);
    \node[labelback,font=\scriptsize] at (1.35,.75) {antihorario};
    \node[labelback,font=\small,text=rust] at (2,-1.05) {terreno (m)};
  \end{scope}
\end{tikzpicture}
\end{center}

\subsection{Afinidades construidas geométricamente}\label{sec:afinidades-construidas}

Una afinidad también puede definirse con regla y escuadra: se fija un eje, una dirección de
desplazamiento y un factor. Los puntos del eje no se mueven; los demás se deslizan en la
dirección elegida.

\begin{definicion}[Afinidad axial]\label{def:06-axial}
Sean \(L\) una recta (el \emph{eje}), \(d\) una dirección no paralela a \(L\) y \(k\neq0\) un
número (la \emph{razón}). Todo punto se escribe de manera única como \(x=a+\beta d\) con
\(a\in L\). La \emph{afinidad axial} de eje \(L\), dirección \(d\) y razón \(k\) es
\[
  f(a+\beta d)=a+k\beta d .
\]
Si \(d\) es perpendicular a \(L\), la afinidad axial es \emph{ortogonal}; si no, es
\emph{oblicua}.
\end{definicion}

\begin{ejemplo}[Una afinidad axial]\label{ej:06-axial}
Tomemos como eje \(L\) el eje \(x\), dirección \(d=(1,1)\) y razón \(k=2\). El punto
\((x,y)\) se descompone como \((x-y,\,0)+y\,(1,1)\), con \(a=(x-y,0)\in L\) y \(\beta=y\). Su
imagen es
\[
  f(x,y)=(x-y,\,0)+2y\,(1,1)=(x+y,\ 2y),
  \qquad
  A=\begin{pmatrix}1&1\\0&2\end{pmatrix},\quad b=0 .
\]
El punto \((1,2)\) tiene \(a=(-1,0)\) y \(\beta=2\); va a \((-1,0)+4(1,1)=(3,4)\). Se desplazó
\((2,2)\), paralelo a \(d\), y su distancia al eje medida en la dirección \(d\) se duplicó. Los
puntos \((0,0)\) y \((3,0)\) del eje quedan fijos, y \(\det A=2\): las áreas se duplican.
\end{ejemplo}

\begin{center}
\begin{tikzpicture}[scale=.9]
  \draw[axisline,-{Latex[length=2mm]}] (0,-.5)--(0,4.6) node[above] {$y$};
  \draw[rust,line width=1.1pt] (-2,0)--(5.4,0) node[right] {eje \(L\)};
  \coordinate (a) at (-1,0);
  \coordinate (p) at (0,0);
  \coordinate (q) at (3,0);
  \coordinate (x) at (1,2);
  \coordinate (xp) at (3,4);
  \draw[auxline] (a)--(3.5,4.5);
  \draw[auxline] (p)--(3.3,4.4);
  \draw[auxline] (p)--(1.4,2.8);
  \draw[auxline] (q)--(.4,2.6);
  \draw[auxline] (q)--(3,4.4);
  \fill[ink] (a) circle (2.2pt) node[below] {$a$};
  \fill[ink] (p) circle (2.2pt) node[below right] {$(0,0)$};
  \fill[ink] (q) circle (2.2pt) node[below] {$(3,0)$};
  \fill[teal] (x) circle (2.5pt) node[labelback,left=2pt] {$x=(1,2)$};
  \fill[rust] (xp) circle (2.5pt) node[labelback,right=2pt] {$f(x)=(3,4)$};
  \draw[mainline,-{Latex[length=2.5mm]}] (x)--(xp)
    node[midway,below right,font=\small] {dirección \(d\)};
\end{tikzpicture}
\end{center}

La recta por \((0,0)\) y \(x\) va a la recta por \((0,0)\) y \(f(x)\); la recta por \((3,0)\) y
\(x\) va a la recta por \((3,0)\) y \(f(x)\). Esta es la construcción gráfica: cada recta corta
al eje en un punto fijo, y su imagen pasa por ese mismo punto.

¿Cuántos datos hacen falta para fijar una afinidad cualquiera? La matriz \(A\) tiene cuatro
números y \(b\) dos: seis incógnitas. Cada punto con su imagen aporta dos ecuaciones, así que
alcanzan tres puntos, siempre que no estén alineados.

\begin{idea}[Idea]
Tres puntos no alineados forman un triángulo, y un triángulo es un sistema de referencia
afín: un vértice hace de origen y los dos lados que salen de él hacen de ejes. Decir adónde
van los tres vértices es decir adónde van el origen y los ejes, y eso fija toda la
afinidad.
\end{idea}

\begin{proposicion}[Una afinidad del plano queda determinada por la imagen de tres puntos no alineados]\label{prop:06-tres-puntos}
Sean \(p_1,p_2,p_3\) no alineados y \(q_1,q_2,q_3\) puntos cualesquiera. Existe una única
transformación \(f(x)=Ax+b\) (con \(A\) una matriz \(2\times2\) cualquiera) tal que
\(f(p_i)=q_i\) para \(i=1,2,3\). Con
\[
  P=\bigl[\,p_2-p_1\ \big|\ p_3-p_1\,\bigr],
  \qquad
  P'=\bigl[\,q_2-q_1\ \big|\ q_3-q_1\,\bigr],
\]
es \(A=P'P^{-1}\) y \(b=q_1-Ap_1\). La afinidad es invertible si y sólo si \(q_1,q_2,q_3\)
tampoco están alineados.
\end{proposicion}

\begin{demostracion}
Como \(p_1,p_2,p_3\) no están alineados, los vectores \(p_2-p_1\) y \(p_3-p_1\) son linealmente
independientes, y la matriz \(P\) que los tiene por columnas es invertible.

\emph{Unicidad.} Si \(f(x)=Ax+b\) cumple \(f(p_i)=q_i\), restando
\(f(p_i)-f(p_1)=A(p_i-p_1)\) se obtiene \(A(p_i-p_1)=q_i-q_1\) para \(i=2,3\); en forma matricial,
\(AP=P'\). Por lo tanto \(A=P'P^{-1}\), y de \(Ap_1+b=q_1\) sale \(b=q_1-Ap_1\). No hay otra
opción.

\emph{Existencia.} Definimos \(A=P'P^{-1}\) y \(b=q_1-Ap_1\). Entonces \(f(p_1)=q_1\), y para
\(i=2,3\), \(f(p_i)=f(p_1)+A(p_i-p_1)=q_1+(q_i-q_1)=q_i\), porque \(AP=P'\) dice exactamente que
\(A\) lleva cada columna de \(P\) a la correspondiente de \(P'\).

\emph{Invertibilidad.} \(\det A=\det P'/\det P\), que es distinto de cero si y sólo si las
columnas de \(P'\) son independientes, es decir, si \(q_1,q_2,q_3\) no están alineados.
\end{demostracion}

\begin{ejemplo}[Afinidad a partir de tres puntos de control]\label{ej:06-tres-puntos}
Una fotografía escaneada se georreferencia con tres puntos de control, cuyas coordenadas de
imagen \((u_i,v_i)\) y de terreno \((E_i,N_i)\) se conocen:
\[
\begin{array}{c|cc|cc}
i & u_i & v_i & E_i & N_i\\ \hline
1 & 10 & 20 & 1005{,}82 & 1990{,}47\\
2 & 110 & 30 & 1056{,}04 & 1989{,}60\\
3 & 40 & 120 & 1024{,}96 & 1942{,}00
\end{array}
\]
Las incógnitas son los seis coeficientes de
\(E=a_{11}u+a_{12}v+b_1\), \(N=a_{21}u+a_{22}v+b_2\). Cada punto aporta una ecuación para \(E\)
y otra para \(N\). Ordenando las incógnitas como \((a_{11},a_{12},b_1,a_{21},a_{22},b_2)\), el
sistema de \(6\times6\) es diagonal por bloques, con el mismo bloque \(U\) dos veces:
\[
  \begin{pmatrix}U&0\\0&U\end{pmatrix}
  \begin{pmatrix}a_{11}\\a_{12}\\b_1\\a_{21}\\a_{22}\\b_2\end{pmatrix}
  =\begin{pmatrix}E_1\\E_2\\E_3\\N_1\\N_2\\N_3\end{pmatrix},
  \qquad
  U=\begin{pmatrix}10&20&1\\110&30&1\\40&120&1\end{pmatrix},
  \qquad \det U=9700 .
\]
El determinante de \(U\) es el doble del área del triángulo de control (en píxeles al
cuadrado); no es nulo porque los puntos no están alineados. Resolviendo
\(U(a_{11},a_{12},b_1)^T=(E_1,E_2,E_3)^T\) y \(U(a_{21},a_{22},b_2)^T=(N_1,N_2,N_3)^T\):
\[
  A=\begin{pmatrix}0{,}498&0{,}042\\0{,}041&-0{,}497\end{pmatrix},
  \qquad
  b=\begin{pmatrix}1000{,}000\\2000{,}000\end{pmatrix},
  \qquad
  \det A=-0{,}249 .
\]
La fórmula de la proposición~\ref{prop:06-tres-puntos} da lo mismo: con
\(P=\begin{pmatrix}100&30\\10&100\end{pmatrix}\) y
\(P'=\begin{pmatrix}50{,}22&19{,}14\\-0{,}87&-48{,}47\end{pmatrix}\), el producto \(P'P^{-1}\)
es la misma \(A\). El píxel mide algo menos de \(0{,}5\ \mathrm{m}\), la imagen está girada
unos grados respecto del Norte y, como en el ejemplo~\ref{ej:06-dron}, \(\det A<0\) por el
eje \(v\) hacia abajo.

\emph{Verificación con un cuarto punto.} El punto \((u,v)=(80,80)\) no se usó para calcular
la afinidad; en el terreno se midió \((1043{,}23,\ 1963{,}49)\). La afinidad predice
\[
  A\begin{pmatrix}80\\80\end{pmatrix}+b
  =\begin{pmatrix}1043{,}20\\1963{,}52\end{pmatrix},
\]
con un residuo (medido menos predicho) de \((0{,}03,\ -0{,}03)\ \mathrm{m}\), de módulo
\(0{,}042\ \mathrm{m}\). Con tres puntos el ajuste es exacto y no deja control; el cuarto
punto es el que permite juzgar la calidad de la georreferenciación.
\end{ejemplo}

\subsection{Rotación alrededor de un centro y escala alrededor de otro}\label{sec:06-rot-escala}

Este caso reúne todo lo anterior: dos afinidades con centros distintos, su composición como
producto de matrices homogéneas y el punto fijo del resultado.

\begin{ejemplo}[Rotación alrededor de un centro y escala alrededor de otro]\label{ej:06-rot-escala}
Se toma una figura en forma de T con vértices
\[
\begin{aligned}
&(-2,2),\ (2,2),\ (2,1),\ (0{,}5,\ 1),\\
&(0{,}5,\ -2),\ (-0{,}5,\ -2),\ (-0{,}5,\ 1),\ (-2,1).
\end{aligned}
\]
Primero se aplica la rotación \(\rhoP{c,45^\circ}\) alrededor de \(c=(3,1)\); después, la
homotecia \(\homot{o}{1/2}\) de centro \(o=(-1,2)\) y razón \(k=\tfrac12\). Con la tabla de la
sección~\ref{sec:puntos-fijos-afin}, sus matrices homogéneas son
\[
M_1=
\begin{pmatrix}
0{,}707&-0{,}707&1{,}586\\
0{,}707&0{,}707&-1{,}828\\
0&0&1
\end{pmatrix},
\qquad
M_2=
\begin{pmatrix}
0{,}5&0&-0{,}5\\
0&0{,}5&1\\
0&0&1
\end{pmatrix}.
\]
Por la proposición~\ref{prop:06-composicion}, la composición \(\homot{o}{1/2}\circ\rhoP{c,45^\circ}\)
tiene matriz
\[
M=M_2M_1=
\begin{pmatrix}
0{,}354&-0{,}354&0{,}293\\
0{,}354&0{,}354&0{,}086\\
0&0&1
\end{pmatrix}.
\]
El orden importa: \(M_1M_2\) tiene la misma parte lineal pero traslación \((0{,}525,\ -1{,}475)\).
El vértice \(q=(2,2)\) sigue
\[
  q=(2,2)\ \longmapsto\ q_1=(1{,}586,\ 1{,}000)\ \longmapsto\ q_2=(0{,}293,\ 1{,}500),
\]
y \(M\) aplicada a \((2,2,1)^T\) da directamente \((0{,}293,\ 1{,}500,\ 1)^T\). La primera
etapa recorre un arco de \(45^\circ\) alrededor de \(c\); la segunda, el segmento de \(q_1\)
a \(o\) hasta su punto medio.

\emph{Punto fijo.} Con \(A=\begin{pmatrix}0{,}354&-0{,}354\\0{,}354&0{,}354\end{pmatrix}\) y
\(b=(0{,}293,\ 0{,}086)\), el sistema \((I-A)p_\ast=b\) da
\[
  p_\ast=\Bigl(1-\tfrac{\sqrt2}{2},\ 1-\tfrac{\sqrt2}{2}\Bigr)\approx(0{,}293,\ 0{,}293).
\]
Ninguna de las dos transformaciones lo deja quieto por separado, pero la composición sí:
\(M\,(0{,}293,\ 0{,}293,\ 1)^T=(0{,}293,\ 0{,}293,\ 1)^T\). Como \(A=\tfrac12R(45^\circ)\), la
composición es una similaridad: gira \(45^\circ\) y reduce a la mitad alrededor de \(p_\ast\).
\end{ejemplo}

\begin{center}
\begin{tikzpicture}[scale=.68]
  \begin{scope}
    \draw[axisline,-{Latex[length=1.8mm]}] (-2.5,0)--(3.6,0);
    \draw[axisline,-{Latex[length=1.8mm]}] (0,-2.5)--(0,2.6);
    \draw[teal,line width=1.3pt]
      (-2,2)--(2,2)--(2,1)--(.5,1)--(.5,-2)--(-.5,-2)--(-.5,1)--(-2,1)--cycle;
    \fill[teal] (2,2) circle (2.5pt) node[labelback,above] {$q$};
    \fill[ink] (3,1) circle (2.5pt) node[labelback,right] {$c$};
    \node[labelback,font=\small] at (.3,-3.0) {figura inicial};
  \end{scope}

  \draw[-{Latex[length=2.5mm]},tealdark,line width=.9pt] (4.0,0)--(5.1,0)
    node[midway,above,font=\small] {\(\rhoP{c,45^\circ}\)};

  \begin{scope}[xshift=6.2cm]
    \draw[axisline,-{Latex[length=1.8mm]}] (-2.5,0)--(3.8,0);
    \draw[axisline,-{Latex[length=1.8mm]}] (0,-4.0)--(0,2.0);
    \draw[tealdark,line width=1.3pt]
      (-1.243,-1.828)--(1.586,1.000)--(2.293,.293)--(1.232,-.768)--
      (3.354,-2.889)--(2.646,-3.596)--(.525,-1.475)--(-.536,-2.536)--cycle;
    \fill[tealdark] (1.586,1.000) circle (2.5pt) node[labelback,above] {$q_1$};
    \fill[ink] (3,1) circle (2.5pt) node[labelback,right] {$c$};
    \node[labelback,font=\small] at (.5,-4.5) {después de rotar};
  \end{scope}

  \draw[-{Latex[length=2.5mm]},tealdark,line width=.9pt] (10.2,0)--(11.3,0)
    node[midway,above,font=\small] {\(\homot{o}{1/2}\)};

  \begin{scope}[xshift=12.4cm]
    \draw[axisline,-{Latex[length=1.8mm]}] (-2.5,0)--(2.0,0);
    \draw[axisline,-{Latex[length=1.8mm]}] (0,-2.6)--(0,2.5);
    \draw[rust,line width=1.3pt]
      (-1.121,.086)--(.293,1.500)--(.646,1.146)--(.116,.616)--
      (1.177,-.445)--(.823,-.798)--(-.237,.263)--(-.768,-.268)--cycle;
    \fill[rust] (.293,1.500) circle (2.5pt) node[labelback,above] {$q_2$};
    \fill[ink] (-1,2) circle (2.5pt) node[labelback,left] {$o$};
    \fill[gold] (.293,.293) circle (2.8pt) node[labelback,right] {$p_\ast$};
    \node[labelback,font=\small] at (-.2,-3.0) {resultado final};
  \end{scope}
\end{tikzpicture}
\end{center}

\section{En profundidad}

Leemos ahora la afinidad por sus autovalores, conectamos los puntos fijos con la matriz
homogénea y ubicamos el capítulo en la escalera \(\R\to\R^2\to\R^n\). Esto prepara el cambio
de base del capítulo~\ref{cap:cambio-base}.

\subsection{La recíproca y el apéndice de composición}\label{sec:06-reciproca}

La proposición~\ref{prop:06-conserva} muestra que toda afinidad conserva rectas y
paralelismo. Vale también la recíproca: una biyección del plano que lleva rectas en rectas
es de la forma \(Ax+b\). Es el teorema fundamental de la geometría afín; su demostración usa
que la única función de \(\R\) en \(\R\) que respeta sumas, productos y orden es la
identidad, y excede el alcance de este apunte. En la práctica, se toma
\(f(x)=Ax+b\) como definición (definición~\ref{def:06-afinidad}) y se verifican las
propiedades geométricas.

El apéndice~\ref{ap:composicion-afin} desarrolla con todo detalle el
ejemplo~\ref{ej:06-rot-escala}: deduce la matriz homogénea de la rotación con centro y la de
la escala con centro, demuestra por bloques que la composición es \(A=kR\),
\(b=k(I-R)c+(1-k)o\), obtiene el punto fijo como \(p_\ast=(I-A)^{-1}b\) y lo identifica con el
autovector de autovalor \(1\) de la matriz homogénea, con los valores de cada paso y su
implementación en planilla de cálculo.

\subsection{La misma idea en cada dimensión}\label{sec:misma-idea}

Casi todo lo que hacemos con matrices en \(\R^2\) y \(\R^3\) ya existe en \(\R\), donde es tan
simple que pasa inadvertido. Si una fórmula se entiende con un solo número, la matriz sólo
agrega direcciones que pueden mezclarse.

{\small
\begin{tabularx}{\textwidth}{>{\bfseries\raggedright\arraybackslash}p{17mm}>{\raggedright\arraybackslash}p{30mm}>{\raggedright\arraybackslash}p{33mm}>{\raggedright\arraybackslash}p{30mm}Y}
\toprule
& \(\R\) & \(\R^2\) & \(\R^3\) y \(\R^n\) & Hecho geométrico\\
\midrule
Afinidad & \(x\mapsto ax+b\): cambio de unidad y de origen en una recta graduada &
\(Ax+b\), \(A\) de \(2\times2\); homogénea \(3\times3\) & \(A\) de \(n\times n\); homogénea
\((n+1)\times(n+1)\) & conserva rectas, paralelas y razones sobre una recta\\
\addlinespace
Identidad & el número \(1\) & la matriz \(I\): \(Iv=v\) & \(I\) de \(n\times n\) & no mueve
nada; neutro de la composición\\
\addlinespace
Inversa & \(1/a\); no existe si \(a=0\) & \(A^{-1}\); no existe si \(\det A=0\) & igual; si
\(\det A=0\) hay una dirección aplastada & deshacer la transformación\\
\addlinespace
Tamaño con signo & \(|a|\) estira; el signo dice si da vuelta & \(\det A\): área orientada
del paralelogramo imagen del cuadrado unidad & volumen orientado & cuánto se agranda una
figura y si cambia la orientación\\
\addlinespace
Punto fijo & \(x_\ast=b/(1-a)\), si \(a\neq1\) & \((I-A)p_\ast=b\) & igual & el lugar donde se
ancla la transformación\\
\addlinespace
Autovector & trivial: toda dirección, con \(\lambda=a\) & sólo algunas direcciones, o
ninguna real si hay giro & una rotación de \(\R^3\) tiene una: el eje & dirección que cae
sobre sí misma\\
\bottomrule
\end{tabularx}
}

% Apéndice: resolución de los ejercicios de afinidades
\section{Resolución de los ejercicios de afinidades}\label{cap:17}\label{ap:ejercicios-afinidades}

\subsection{Ejercicio 1: tres reflexiones puntuales}

\textbf{Enunciado.} Dados tres puntos no alineados \(a,b,c\), hallar el punto fijo de
\(\sigma_a\circ\sigma_b\circ\sigma_c\) y verificar que forma un paralelogramo con ellos.

\textbf{Solución algebraica.} La reflexión puntual de centro \(p\) es
\[
  \sigma_p(x)=2p-x.
\]
Entonces:
\[
\begin{aligned}
\sigma_a\sigma_b\sigma_c(x)
&=\sigma_a\!\left(\sigma_b(2c-x)\right)\\
&=\sigma_a(2b-2c+x)\\
&=2a-2b+2c-x\\
&=2(a-b+c)-x.
\end{aligned}
\]
La composición es una reflexión puntual de centro
\[
  \boxed{d=a-b+c}.
\]
Su punto fijo es \(d\), porque \(\sigma_d(d)=d\). Además:
\[
  b+d=a+c.
\]
Las diagonales \(ac\) y \(bd\) tienen el mismo punto medio, de modo que \(a,b,c,d\) forman
un paralelogramo.

\textbf{Interpretación geométrica.} Tomando \(b\) como origen provisional:
\[
  \overrightarrow{bd}
  =\overrightarrow{ba}+\overrightarrow{bc}.
\]
El punto \(d\) se obtiene mediante la regla del paralelogramo para sumar los vectores que van
desde \(b\) hacia \(a\) y hacia \(c\).

\subsection{Ejercicio 2: dos reflexiones puntuales}

\textbf{Enunciado.} Dados \(a\neq b\), demostrar que
\(\sigma_a\circ\sigma_b\) no tiene puntos fijos.

\textbf{Solución.}
\[
\sigma_a\sigma_b(x)=x+2(a-b).
\]
Si existiera un punto fijo:
\[
x+2(a-b)=x
\quad\Longrightarrow\quad
a=b.
\]
Esto contradice \(a\neq b\). Por lo tanto:
\[
  \boxed{\operatorname{Fix}(\sigma_a\sigma_b)=\varnothing
  \quad\text{si }a\neq b.}
\]

\textbf{Interpretación geométrica.} La composición es una traslación no nula. Todos los puntos
se desplazan por el mismo segmento orientado, por lo que ninguno puede permanecer en su lugar.
Si \(a=b\), las dos reflexiones se anulan y la composición es la identidad.

\subsection{Ejercicio 3: cuatro reflexiones en rectas}

\begin{atencion}[Ambigüedad del enunciado]
El texto comienza diciendo ``dadas cuatro rectas \(a,b,c,d\)'', pero termina llamando a
\(\varphi_i\) ``reflexión alrededor del punto \(i\)''. Por el contexto y la notación
\(\varphi\), se interpreta que son reflexiones respecto de las cuatro rectas.
\end{atencion}

Cada reflexión en una recta es una afinidad
\[
  \varphi_i(x)=H_ix+q_i,
\]
donde \(H_i\) es una matriz ortogonal de determinante \(-1\). La composición
\[
  F=\varphi_a\varphi_b\varphi_c\varphi_d
\]
tiene la forma
\[
  F(x)=Ax+t,
\]
con
\[
  A=H_aH_bH_cH_d,
\]
\[
  t=H_aH_bH_cq_d+H_aH_bq_c+H_aq_b+q_a.
\]
Como intervienen cuatro reflexiones,
\[
  \det A=(-1)^4=1.
\]
La transformación resultante preserva orientación y es una rotación o una traslación.

\textbf{Método analítico.} Los puntos fijos satisfacen
\[
  (I-A)x=t.
\]
\begin{itemize}
  \item Si \(I-A\) es invertible, existe un único punto fijo:
  \[
    \boxed{x_\ast=(I-A)^{-1}t}.
  \]
  La composición es una rotación con centro \(x_\ast\).
  \item Si \(A=I\) y \(t\neq0\), la composición es una traslación y no tiene punto fijo.
  \item Si \(A=I\) y \(t=0\), es la identidad y todos los puntos son fijos.
\end{itemize}

\textbf{Método gráfico.} Elegir dos puntos \(p,q\), construir sucesivamente sus cuatro
reflejos hasta obtener \(F(p)\) y \(F(q)\), y trazar las mediatrices de
\(\overline{pF(p)}\) y \(\overline{qF(q)}\). Si se cortan, la intersección es el centro de
rotación. Si son paralelas, el resultado es una traslación y no existe centro finito.

\subsection{Ejercicio 4: autovalor de una rotación en \(\R^3\)}

Sea \(R\in SO(3)\). Al ser una matriz real de dimensión impar, su polinomio característico
tiene al menos una raíz real. Como \(R\) es ortogonal, todo autovalor tiene módulo \(1\), por
lo que el autovalor real debe ser \(1\) o \(-1\).

Los otros dos autovalores, si no son reales, forman un par conjugado
\(e^{i\theta},e^{-i\theta}\), cuyo producto es \(1\). Como
\[
  \det R=1,
\]
el autovalor real restante debe ser
\[
  \boxed{\lambda=1}.
\]
Su autovector es la dirección del eje de rotación.

\begin{atencion}[Sobre la palabra ``único'']
Para una rotación no trivial de ángulo distinto de \(0\), \(1\) es el único autovalor real
asociado al eje, salvo casos degenerados. Para \(R=I\), todos los vectores son propios con
autovalor \(1\); por eso el enunciado necesita excluir la identidad si pretende unicidad del
eje.
\end{atencion}

\subsection{Ejercicio 5: matriz ortogonal que invierte orientación}

El enunciado afirma que una matriz ortogonal de \(\R^3\) con determinante \(-1\) tiene tres
autovalores reales iguales a \(1\) o \(-1\). Esa afirmación no es cierta en general.

Por ser real y de dimensión impar, la matriz posee al menos un autovalor real. Por ser
ortogonal, sus autovalores tienen módulo \(1\), y por tener determinante \(-1\), siempre
existe al menos un autovalor
\[
  \boxed{\lambda=-1}.
\]
Los otros dos pueden ser
\[
  e^{i\theta},\qquad e^{-i\theta},
\]
que no son reales cuando \(\theta\notin\{0,\pi\}\). Geométricamente, esto describe una
\textbf{rotorreflexión}: una rotación en un plano combinada con una inversión en la dirección
perpendicular.

\textbf{Contraejemplo.}
\[
Q=
\begin{pmatrix}
\cos\theta&-\sin\theta&0\\
\sin\theta& \cos\theta&0\\
0&0&-1
\end{pmatrix}
\]
es ortogonal y \(\det Q=-1\), pero sus autovalores son
\[
  e^{i\theta},\quad e^{-i\theta},\quad -1.
\]
Solo para casos especiales, como una reflexión en un plano, los tres son reales:
\[
  1,\ 1,\ -1.
\]

\subsection{Ejercicio 6: recuperar una afinidad desde tres puntos}

Sean \(a,b,c\) tres puntos no alineados y sean \(a',b',c'\) sus imágenes. Se busca
\[
  \gamma(x)=Ax+t.
\]
Restando la ecuación de \(a\) a las de \(b\) y \(c\):
\[
  b'-a'=A(b-a),
  \qquad
  c'-a'=A(c-a).
\]
Construimos las matrices cuyas columnas son los lados de ambos triángulos:
\[
  P=
  \begin{pmatrix}\vert&\vert\\b-a&c-a\\\vert&\vert\end{pmatrix},
  \qquad
  P'=
  \begin{pmatrix}\vert&\vert\\b'-a'&c'-a'\\\vert&\vert\end{pmatrix}.
\]
Como \(a,b,c\) no están alineados, \(P\) es invertible. Entonces:
\[
  \boxed{A=P'P^{-1}},
\]
y usando la correspondencia \(a\mapsto a'\):
\[
  \boxed{t=a'-Aa}.
\]
La transformación queda completamente determinada:
\[
  \boxed{\gamma(x)=P'P^{-1}x+a'-P'P^{-1}a}.
\]

\textbf{Comprobación de ausencia de puntos fijos.} Se intenta resolver
\[
  (I-A)x=t.
\]
Que la afinidad no tenga puntos fijos significa que este sistema es incompatible. En
coordenadas homogéneas:
\[
\begin{pmatrix}x'\\y'\\1\end{pmatrix}
=
\begin{pmatrix}
A_{11}&A_{12}&t_x\\
A_{21}&A_{22}&t_y\\
0&0&1
\end{pmatrix}
\begin{pmatrix}x\\y\\1\end{pmatrix}.
\]

\begin{resultado}[Estado de la guía de afinidades]
Los seis ejercicios quedan resueltos de forma algebraica y geométrica. Los ejercicios 3, 4
y 5 incluyen además las salvedades necesarias para interpretar correctamente sus enunciados.
\end{resultado}

\section{Ejercicios con solución}

En cada ejercicio conviene dibujar los puntos iniciales y finales, indicar el orden de
composición y comprobar el resultado numérico con la matriz homogénea.

\begin{ejercicio}
Escribir en coordenadas homogéneas la rotación de \(90^\circ\) alrededor de \(c=(1,2)\)
seguida de la traslación de vector \((3,0)\). Calcular la imagen de \((0,0)\) y de \(c\), y el
punto fijo de la composición.
\solucion{\(M=\begin{pmatrix}0&-1&6\\1&0&1\\0&0&1\end{pmatrix}\) (la rotación sola tiene
última columna \((3,1)\)); \((0,0)\mapsto(6,1)\), \(c\mapsto(4,2)\); punto fijo
\((2{,}5,\ 3{,}5)\): la composición es una rotación de \(90^\circ\) con otro centro.}
\end{ejercicio}

\begin{ejercicio}
Hallar la afinidad que lleva \((0,0)\mapsto(1,1)\), \((1,0)\mapsto(3,2)\) y
\((0,1)\mapsto(0,3)\).
\solucion{\(A=\begin{pmatrix}2&-1\\1&2\end{pmatrix}\), \(b=(1,1)\). Es una similaridad: gira
y escala por \(\sqrt5\).}
\end{ejercicio}

\begin{ejercicio}
Hallar el punto fijo de \(f(x)=\begin{pmatrix}2&1\\0&3\end{pmatrix}x+\begin{pmatrix}1\\-4\end{pmatrix}\)
y verificarlo.
\solucion{\((I-A)p=b\) da \(p_\ast=(-3,2)\); en efecto, \(Ap_\ast+b=(-4,6)+(1,-4)=(-3,2)\).}
\end{ejercicio}

\begin{ejercicio}
Para \(A=\begin{pmatrix}1&2\\3&4\end{pmatrix}\), calcular \(\det A\), transformar el triángulo
\((0,0),(1,0),(0,1)\) y comparar áreas y orientación.
\solucion{\(\det A=-2\). La imagen es \((0,0),(1,3),(2,4)\); el área pasa de \(0{,}5\) a \(1\)
y la orientación se invierte (área orientada \(-1\)).}
\end{ejercicio}

\begin{ejercicio}
Hallar la inversa de \(f(x)=\begin{pmatrix}2&0\\1&1\end{pmatrix}x+\begin{pmatrix}4\\-1\end{pmatrix}\),
en forma afín y en forma homogénea.
\solucion{\(f^{-1}(y)=A^{-1}(y-b)\) con \(A^{-1}=\begin{pmatrix}0{,}5&0\\-0{,}5&1\end{pmatrix}\)
y traslación \(-A^{-1}b=(-2,3)\); \(M^{-1}=\begin{pmatrix}0{,}5&0&-2\\-0{,}5&1&3\\0&0&1\end{pmatrix}\).}
\end{ejercicio}

\begin{ejercicio}
Hallar la afinidad que lleva el triángulo \((0,0),(2,0),(0,1)\) al triángulo
\((1,1),(5,2),(2,4)\), usando \(A=P'P^{-1}\).
\solucion{\(P=\begin{pmatrix}2&0\\0&1\end{pmatrix}\), \(P'=\begin{pmatrix}4&1\\1&3\end{pmatrix}\),
\(A=\begin{pmatrix}2&1\\0{,}5&3\end{pmatrix}\), \(b=(1,1)\); \(\det A=5{,}5\).}
\end{ejercicio}
